\section{FTP file-transfer benchmarking}\label{sec:ftp-benchmark}

\subsection{Inputs and experiment method}
The saved performance experiment was run on 15 November 2026 using the terminal benchmark driver and real loopback TCP transfers~\cite{ftpmeasure}. The source files are playable media generated offline with FFmpeg: a stereo 48 kHz, 16-bit PCM WAV signal; a moving 640$\times$360 raw YUV420 video in an AVI container; and an H.264 MP4 counterpart of the video. They contain generated signals/patterns rather than recordings, and are not small files padded with arbitrary bytes.

\begin{table}[H]
\centering\small
\caption{Actual input sizes. One MiB is $2^{20}=1,048,576$ bytes.}
\begin{tabularx}{\linewidth}{@{}l r r Y@{}}
\toprule
Input & Bytes & MiB & Role \\
\midrule
\code{audio.wav} & 134,217,806 & 128.000 & Large uncompressed audio. \\
\code{video.avi} & 134,454,314 & 128.226 & Large raw moving video. \\
\code{video.mp4} & 1,831,488 & 1.747 & Already-compressed control; genuinely smaller. \\
\bottomrule
\end{tabularx}
\end{table}

There are two repetitions for each input. Each comparison tests the original source, then gzip level~6: $3\ \text{files}\times2\ \text{repetitions}\times2\ \text{modes}=12$ measured runs. For every run the engine uploads a unique temporary filename through the implemented client, downloads it again, decompresses when necessary, and compares both the original byte length and SHA-256 digest. All 12 saved runs passed verification. Remote verification payloads and decompressed data are bounded to their known expected lengths.

Gzip is preprocessing performed by the benchmark. It produces an ordinary compressed file that FTP transfers in binary mode; it is not a new FTP wire command or an implementation of MODE~Z. Temporary benchmark remote files are deleted after use, while source files remain unchanged.

\begin{algorithm}[H]\small
\caption{One original/gzip comparison}
\begin{algorithmic}[1]
\Require source file and a factory returning an authenticated FTP client
\State stream the original source to compute expected size and SHA-256
\For{mode in [plain, gzip]}
 \State start per-mode timer; prepare payload (gzip level 6 when selected)
 \State connect/login; record preparation time
 \State upload a unique remote filename, sampling payload progress
 \State record FTP upload time and preparation-plus-upload time
 \State download the same payload with an expected-length bound
 \State decompress if selected; verify original length and SHA-256
 \State record verification timings and remove this run's remote file
\EndFor
\State export run rows, progress samples and figures
\end{algorithmic}
\end{algorithm}

\subsection{Metric definitions}
Let $S$ be the original source bytes, $P$ the FTP payload bytes, $t_p$ preparation time, and $t_u$ the measured FTP upload-call duration. The report distinguishes two rates:
\[
 R_{\mathrm{FTP}}=\frac{P}{2^{20}t_u}\ \mathrm{MiB/s},\qquad
 R_{\mathrm{effective}}=\frac{S}{2^{20}t_{\mathrm{total}}}\ \mathrm{MiB/s},\qquad
 \mathrm{saving}=100\left(1-\frac{P}{S}\right)\%.
\]
Here $t_{\mathrm{total}}$ runs from per-mode preparation through upload completion. It includes gzip work and connection/login setup before the upload; $t_u$ includes TYPE, data negotiation, payload transfer and the final completion reply inside the client's upload call. Initial source hashing, verification download, verification hashing/decompression, and cleanup are recorded separately rather than hidden inside the upload rate. $t_{\mathrm{total}}$ is approximately $t_p+t_u$, with small intervening measurement overhead.

\begin{table}[H]
\centering\small
\caption{Exported measurements and what they cover.}
\begin{tabularx}{\linewidth}{@{}l Y@{}}
\toprule
Field & Meaning \\
\midrule
\code{compression\_seconds} & Gzip work only; zero for the original file. \\
\code{preparation\_seconds} & Per-mode preprocessing and FTP connection/login. \\
\code{transfer\_seconds} & The FTP upload call, through successful completion. \\
\code{total\_seconds} & Preparation through upload completion. \\
\code{download\_seconds} & Verification retrieval through FTP. \\
\code{verification\_seconds} & Streamed size/hash check, including decompression if needed. \\
\code{round\_trip\_seconds} & Per-mode preparation, upload, retrieval and content verification. \\
\bottomrule
\end{tabularx}
\end{table}

\subsection{Measured results and graphs}
Tables~\ref{tab:ftp-compression} and~\ref{tab:ftp-timing} show arithmetic means across the two repetitions. The original rows and sampled curves are retained in \code{data/runs.csv}, \code{data/samples.csv}, and \code{data/results.json}; the media manifest records codecs, sizes, generation commands and hashes.
Manifest file paths are normalized to the portable example directory. Measured byte counts, hashes, codec information and generation timings are unchanged.

\begin{table}[H]
\centering\small
\caption{Mean payload sizes and reduction over two repetitions.}
\label{tab:ftp-compression}
\begin{tabularx}{\linewidth}{@{}Y r r r@{}}
\toprule
Source & Original (MiB) & Gzip (MiB) & Bytes saved \\
\midrule
\code{audio.wav} & 128.000 & 3.751 & 97.070\% \\
\code{video.avi} & 128.226 & 6.118 & 95.228\% \\
\code{video.mp4} & 1.747 & 1.744 & 0.173\% \\
\bottomrule
\end{tabularx}
\end{table}
\begin{table}[H]
\centering\small
\caption{Mean FTP payload rate and preparation-plus-upload time. Rates are arithmetic means of run rates.}
\label{tab:ftp-timing}
\begin{tabularx}{\linewidth}{@{}Y r r r r@{}}
\toprule
& \multicolumn{2}{c}{FTP payload rate (MiB/s)} & \multicolumn{2}{c}{Preparation + upload (s)} \\
\cmidrule(lr){2-3}\cmidrule(l){4-5}
Source & Original & Gzip & Original & Gzip \\
\midrule
\code{audio.wav} & 1962.33 & 90.70 & 0.0689 & 0.8808 \\
\code{video.avi} & 1841.61 & 147.23 & 0.0710 & 1.1468 \\
\code{video.mp4} & 42.72 & 42.04 & 0.0413 & 0.0900 \\
\bottomrule
\end{tabularx}
\end{table}


\begin{figure}[H]
\centering
\includegraphics[width=\linewidth]{ftp_benchmark_compression.png}
\caption{Original versus gzip payload sizes for each saved file/repetition. The generated WAV and raw AVI patterns compress substantially; the MP4 counterpart changes very little. Source: measured experiment rows, 15 November 2026.}
\label{fig:ftp-compression-graph}
\end{figure}

\begin{figure}[H]
\centering
\includegraphics[width=\linewidth]{ftp_benchmark_throughput.png}
\caption{Measured FTP payload throughput and effective original-byte rate. A smaller gzip payload can have a different payload rate and a lower effective rate once compression time is counted. These are cached loopback measurements, not physical-link capacity.}
\label{fig:ftp-throughput-graph}
\end{figure}

\begin{figure}[H]
\centering
\includegraphics[width=\linewidth]{ftp_benchmark_durations.png}
\caption{Preparation, upload, verification download and content-verification durations. Keeping phases separate exposes the gzip preprocessing cost instead of attributing it to the network.}
\label{fig:ftp-duration-graph}
\end{figure}

\begin{figure}[H]
\centering
\includegraphics[width=\linewidth,height=0.77\textheight,keepaspectratio]{ftp_benchmark_time-series.png}
\caption{Cumulative FTP upload progress over elapsed upload time. Each original/gzip trace starts at its own time zero; the vertical axis is payload MiB, not MiB/s. The slope reflects progress through the local transfer, and the last sample includes completion-reply latency.}
\label{fig:ftp-timeseries-graph}
\end{figure}

The WAV payload fell to approximately 3.751~MiB (97.070\% smaller) and raw AVI to 6.118~MiB (95.228\% smaller). The already-compressed MP4 saved only 0.173\%. In every case gzip took longer from preparation to upload completion: about 0.881~s versus 0.069~s for WAV, 1.147~s versus 0.071~s for AVI, and 0.090~s versus 0.041~s for MP4. Compression saved bytes but did not save total upload-path time on this machine's local path.

\subsection{Interpretation and limitations of the measurements}
Large-file plain uploads ranged roughly from 1,551 to 2,373~MiB/s. Such high values reflect the operating system's loopback path and warm filesystem caches; they do not establish Ethernet/Wi-Fi capacity or disk durability. There were only two repetitions, plain always preceded gzip, and background load/cache state were not controlled as a formal statistical study. The synthesized signals also contain regular structure, so their compression ratios cannot be generalized to all audio/video.

The benchmark measures individual transfers, not concurrent-user saturation. It does not include a physical second-machine LAN trial, TCP/IP header accounting, or a claim that a file was forced to stable storage. Its useful conclusions are narrower: the implementation streamed real media through FTP, all round trips verified, raw patterned media compressed well, and the measured gzip CPU/setup cost outweighed loopback transfer savings.

\subsection{Reproducing the experiment from the terminal}
Run from the FTP directory, using fresh generated-media and generated-results directories. FFmpeg and FFprobe are needed only for fixture generation. Matplotlib is needed only for \code{--plots}; CSV/JSON export uses the standard library.
\begin{lstlisting}[style=ftpterm]
cd Group_Project/FTP
python3 tools/media_fixtures.py \
  --output benchmarks/generated-media --mib 128 --mp4
python3 -m ftp_lab.benchmark --local \
  benchmarks/generated-media/audio.wav \
  benchmarks/generated-media/video.avi \
  benchmarks/generated-media/video.mp4 \
  --repeat 2 --output benchmarks/generated-results --plots
\end{lstlisting}
For an actual second computer's FTP server, omit \code{--local}, supply \code{--host}, \code{--port}, \code{--user}, and optionally \code{--remote-dir}, and provide its lab password through \code{FTP\_PASSWORD} or a prompt. The target folder must permit upload, retrieval and cleanup. Keep those new observations separate from the saved loopback results.
